\section{Randomized feedback and concurrent execution study}\label{supp:randomized-feedback}
The paired design tests how the information forwarded after a tool response changes the model-link observation. It also measures the consequences of sharing connections and adding model-link shaping. The synthetic business actions execute in isolated fixtures; the tools are served by the MCP SDK and the workflow is executed by LangGraph. The earlier serial padding, common-prefix and fixed-composition studies are indexed in the artifact archive; their populations are not combined with this study.

\subsection{Task population and assignment}
Six actions are available: refund, address update, escalation, replacement shipment, cancellation and receipt. Each task has one to three prerequisites sampled from account, policy, approval and status checks. A task assigned a temporary check error requires a single retry after receiving that terminal response. A correct business action changes the corresponding fixture field. The evaluator records the requested action, applied action, prerequisite completion and final state separately.

The local cohort contains 312 task instances in 64 sessions per model and condition. The training, validation and test partitions contain 32, 8 and 24 sessions, respectively, with 156, 36 and 120 tasks. Each class has 26 training, six validation and 20 test instances. Class counts vary within sessions: the complete cohort contains ten empty sessions, 41 sessions with repeated classes and 59 missing at least one class. Four test sessions are empty, leaving 20 nonempty clusters for supplied-window test accuracy. The test partition uses four held-out request phrasings per action. Sessions vary in arrivals and contents but do not constitute independent application domains.

Arrivals follow the frozen session schedule generated with exponential interarrival times at rate .55/s. Idle periods of 3--7 seconds precede and follow the work; an empty session lasts 12 seconds. Independent MCP \texttt{tools/list} background requests follow a .35/s exponential process. These are live SDK requests on the captured service, not a sample of enterprise background traffic. A session has 0--10 task arrivals; the concurrent conditions use up to four workers. Condition order rotates by session; the local models run as separate cohorts. Empty sessions remain in duration and false-alarm denominators.

\subsection{Feedback isolation and intervention}
The initial instruction, system message and admissible initial action are identical across tasks. The fixture becomes available only after the initial model response terminates. The first tool response then contains the customer request, required checks, order context and three to seven supplementary records. The records use service-specific schemas with randomized content; they deliberately provide a recoverable dependency to test, rather than estimate the frequency of such schemas in deployed services.

\emph{Full} forwards that response to the next model request. \emph{Projected} applies a fixed allowlist after receiving the same complete tool response. It retains \texttt{customer\_request}, \texttt{required\_checks}, order reference, amount and address, plus check/result fields, while removing \texttt{supporting\_records}. Task seeds and tool fixtures are paired across conditions. The audit checks equality of fixture hashes, equality of unprojected fields, and the exact allowlist operation. Projection preserves the business information needed by the designed task; measured completion determines whether the model still uses it successfully.

Both serial conditions pad tool requests to a 4,096-byte quantum and responses to a 16,384-byte quantum. JSON whitespace and SSE-comment padding are removed by parsing before the structured tool result is forwarded. \emph{Dedicated} uses separate connections for each worker and service; \emph{Shared} uses a single connection per service. \emph{Joint} adds model-request padding to a 32,768-byte quantum and emits model SSE frames of 512 bytes every 25 ms for at least 40 frames. Responses extending beyond the minimum schedule remain extended and are counted as overruns. Padding quanta can be exceeded. Thus the intervention is an implemented shaping policy with measured residual signals, not a constant-time channel.

\subsection{Execution and transport}
LangGraph 1.0.3 controls alternating model/tool nodes. MCP SDK 1.21.1 serves Streamable HTTP through Hypercorn 0.17.3. Local model requests go through a TLS~1.3/HTTP/2 gateway that forwards incremental Ollama output as SSE. Local model identities and weights are fixed for Qwen3:8B and Llama3.1:8B. The model chooses among unfinished prerequisites and then all six business actions; the host guard enforces protocol prerequisites without supplying the correct action label. The eight-step cap and functional errors remain part of the outcome population.

Readiness tasks are excluded from the formal partitions. The initial readiness run exposed premature commitments and repeated ticket retrieval; the documented amendment added state-dependent admissible-action constraints before the formal cohort was frozen. A second readiness run checks that the workflow and both models can complete the tasks. Readiness outcomes are preserved independently of formal attack scores.

Capture retains encrypted TCP payload size, relative timestamp, direction and physical connection grouping. Packet features never use a decrypted HTTP/2 stream identifier, task-to-stream mapping, tool plaintext or requested class. A separate private wire check uses TLS keys to verify TLS~1.3, HTTP/2, overlapping streams and incremental delivery. Those verification inputs are excluded from the attack representation.

The formal local capture contains 811,401 payload-bearing TCP packets across 640 sessions. Request and state checks account for all 3,120 executions and 13,680 terminal model requests. Wire verification finds 2,048 TLS~1.3 negotiations selecting HTTP/2 and 39,854 closed HTTP/2 request streams. All model responses contain at least two nonempty DATA frames with a positive first-to-last span. In each shared condition, 107 of 128 sessions contain overlapping model requests on one connection, reaching four active streams. Serial and dedicated conditions have at most one active model request per connection. These checks establish the implemented protocol and overlap; the capture does not provide a packet-drop counter.

\subsection{Recognition models and selection}
Supplied-window attribution has three units: the initial model call, the first model call after the ticket, and the complete execution. Prefix and post-ticket attribution use model-link metadata; complete execution considers the tool link, model link and joint view. A window includes all observed connections in its interval, so concurrent observations include other active tasks. Phase attribution is evaluated for serial Full/Projected, as specified by the mechanism contrasts. Shared/Joint are evaluated through complete windows and continuous localization. Their post-ticket timestamps exist for ground truth, but they do not identify a transit-visible per-task stream: a phase score would mix concurrent tasks. No Joint post-ticket outcome is claimed. Missing phases in declared comparisons remain represented instead of being excluded.

The primary models are standardized logistic regression (LR) and the released 175-feature k-FP extractor~\cite{hayes2016kfp} with a 1,000-tree random forest. LR selects $C\in\{.1,1,10\}$ by validation accuracy, breaking ties toward smaller $C$. The k-FP implementation is version-pinned; Python iterator compatibility changes leave its features unchanged. Each encrypted TCP-payload packet supplies a signed event and relative time to that extractor. Traces with fewer than three packets, only one direction or zero duration receive a flagged zero feature vector. Joint k-FP features concatenate per-link vectors. The closed-world forest evaluation is distinct from k-FP's open-world decision procedure.

Additional attacks are a 500-tree forest on 33 aggregate features and a temporal CNN. The CNN receives at most 512 ordered packets (first 384 and final 128 when truncated), represented by signed log size, log interarrival time, log elapsed time, service role and validity mask. Three convolutional layers with kernel width five precede pooled and fully connected layers. Channel width is selected from 32 and 64; Adam uses learning rate .001, weight decay .001 and batch size 32 for up to 40 epochs. Validation accuracy, then loss and earlier epoch, determines the checkpoint. This is a temporal baseline, not a reproduction of DF or Var-CNN. Its low held-out scores do not establish training failure; the saved selections and outcomes remain reported with the other attacks.

All validation selections are saved before any test score is computed. Each condition is fitted separately; a second evaluation transfers the fitted Full predictor unchanged. Both results are retained. The majority baseline is determined by training labels. Test accuracy, macro-F1, complete confusion matrices and paired correctness are released for every declared learner and view.

\subsection{Continuous localization and uncertainty}
Continuous detectors receive session boundaries and connection identities but no task boundaries or private stream assignment. A ground-truth execution interval starts at worker admission and ends at its terminal outcome; a scheduled task waiting for a worker has not yet entered that interval. Six binary LR models and a multilabel forest use .25-second grids with 1-, 2- or 4-second lookback. Validation event F1 at intersection-over-union (IoU) .5 selects the lookback and threshold, with ties favoring a higher threshold and shorter lookback. Threshold candidates span .1--.9. Gap detection with a supplied-window classifier and constant class-presence outputs provide additional controls.

Predicted events are matched one-to-one to true task instances using maximum-cardinality matching subject to the IoU threshold, with overlap as a tie-break. Duplicate same-class instances remain distinct truths. Binary class-activity models can merge overlapping instances; the missed-instance penalties are retained. We report IoU .3/.5/.7, precision, recall, F1, mean false events per session, duration-normalized false events per hour, matched endpoint errors and empty-session alarms. A strong supplied-window score is not substituted for this event-level evaluation.

Paired uncertainty resamples complete test sessions 2,000 times. LR additionally uses 300 paired training/test session resamples and refits the predictor, keeping its validation-selected hyperparameter fixed. These intervals cover the stated resampling populations, conditional on the designed workload and selected model family. The two feedback contrasts, two primary attacks and two local execution models form one eight-comparison Holm family. Session-level paired sign permutations use 10,000 draws on saved predictions, so their $p$-values condition on the fitted attacks; within-session label permutations use 2,000 draws. A ten-percentage-point contrast was declared practically relevant before scoring. RF/CNN and deployment comparisons remain exploratory; an interval crossing zero does not establish equivalence.

A post-scoring numerical audit enumerates the same session-sign distribution exactly for every declared primary contrast. It retains the original Monte Carlo values, predictions, clusters and multiplicity families. For the smallest local adjusted value, the approximation is .049595 and the exact value is .050781. This threshold sensitivity precludes treating that comparison as a robust rejection. The artifact provides integer numerators, denominators and an independently checked enumeration routine.

Resource accounting separates active execution from session occupancy. Task time runs from worker admission to completion and excludes waiting for an available worker. Session time spans connection setup, arrivals, execution, background traffic and idle margins; its sum divided by the task count gives an amortized session cost, not a per-task response latency. Each captured packet is charged once to its session, including empty sessions. Paired session resampling preserves concurrent workloads without summing overlapping task windows.

\subsection{Hosted retest and review boundary}
The bounded remote plan uses the same workflow and paired Full/Projected intervention with an authorized account and synthetic fixtures. Its separate 96-task population has 36 training, 12 validation and 48 test tasks per condition across 30 sessions. Eighteen test sessions include two empty sessions, leaving 16 nonempty test clusters. Training and validation sessions contain the remaining 36 and 12 tasks. The hosted model identifier is a mutable alias. Fixture seeds remain paired, but provider requests do not enforce a generation seed. Capture observes the local tunnel path plus the local tool service. This retest changes the backend and network together. Provider research permission and terms applicability remain unverified. A missing response terminal is retained as unknown and stops collection without automatic replay.

All 192 formal remote attempts have terminal outcomes, accounting for 800 model requests. Five invalid model outputs terminate their tasks and remain in the intended-label denominator. Full completes 92 of 96 goals and Projected 95 of 96; the held-out counts are 46/48 and 47/48. Wire checks verify TLS~1.3/HTTP/2 on the captured services, 112 negotiations selecting HTTP/2 and 3,196 closed request streams. Every model response has multiple nonempty DATA frames with distinct first and last timestamps. These checks concern the observed protocol, not a fixed provider model or physical-NIC timing guarantee.

The remote test schedule was amended before remote requests or formal attack scoring. The same 48 test instances were redistributed from seven to sixteen nonempty sessions, retaining their labels, fixtures, seeds and partition. With seven clusters, the smallest exact two-sided sign-test value is $2/2^7$, whose first four-comparison Holm step exceeds .05. Sixteen clusters remove that resolution constraint; they do not guarantee power for a ten-point effect. The artifact preserves both plans and the amendment. Local collection, attack selection, multiplicity and the remote request budget remain fixed.

The offline review package includes payload-free relative packet metadata, designed cohort specifications, features, predictions and outcome records. Credentials, TLS keys, raw captures and transcripts remain private. Private wire receipts document checks on excluded inputs; public metadata checks provide the reproducible analysis path. Collection and analysis snapshots, hashes and implementation amendments are stored with the artifact rather than embedded in the scientific narrative.

\subsection{Verified outcomes and all primary comparisons}
The formal cohorts contain 3,120 local and 192 remote attempts, all with terminal records. Local business-goal completion is 1,555/1,560 for Qwen3:8B and 1,542/1,560 for Llama3.1:8B. Remote completion is 187/192. Wrong actions and invalid model outputs remain in the inference population. Table~\ref{tab:randomized-costs} separates all-run completion from held-out completion and reports service time, session time and bytes.

Table~\ref{tab:randomized-primary} contains every primary phase and forwarding comparison. The local LR Full--Projected point differences are 16.7 and 15.0 percentage points, with unadjusted training/test-refit intervals [4.0,29.8] and [3.1,32.8] points. All exact Holm-adjusted primary values exceed .05. The remote LR difference is 10.4 points with interval [-3.2,41.5]; its k-FP difference reverses direction. The intervals and permutation tests condition on different fitted/resampled quantities; the unadjusted intervals do not override the complete multiplicity-adjusted tests.

Table~\ref{tab:randomized-sign-precision} retains the original approximation beside exact enumeration. Table~\ref{tab:randomized-all-recognition} gives every matched and frozen recognition outcome, including the exploratory random forest and CNN. Table~\ref{tab:phase-complete} places complete-execution outcomes beside prefix and first-feedback outcomes. Model-only LR changes from 46/120 to 44/120 for Qwen, 62/120 to 33/120 for Llama, and 9/48 to 9/48 for the hosted model. Complete joint-view counts are 54 to 39, 53 to 29 and 12 to 6, respectively. These views include later requests, responses and decisions. Their heterogeneity rules out characterizing the whole intervention using only its first-feedback point difference. The model's responses and retained task fields are possible sources of residual recognition, but the current bidirectional representation does not attribute that recognition to either one.

In the local Shared joint view, LR and forest event detectors obtain 35.0--43.3\% recall with 4.3--18.6\% precision at IoU .5. The complete event table includes each condition, permitted view, baseline and empty-session count. Mean false events per session divides the aggregate false count by all 24 local test sessions, including four empty sessions (18 and two remotely). Hourly rates remain normalized by the corresponding recorded duration. Session lengths and task loads define these denominators; neither unit supplies enterprise alarm performance. A deployment-prior adjustment would require an event-arrival model and false-alarm process, rather than substituting a classification prior into these event metrics. Table~\ref{tab:randomized-events} places IoU .3/.5/.7 side by side for shared and jointly shaped connections. At .3, the all-active Shared F1 (.295/.508 for Qwen/Llama) exceeds both learned detectors; at .5, the ordering reverses. The same predictions are scored at every threshold. The artifact contains all views, per-session predictions and intervals. A loose threshold can reward coarse occupancy even with random task composition; the .5 primary endpoint remains unchanged.

Projection's local model-byte ratio is approximately .78. Joint shaping increases model-link bytes by 19.5 and 22.5 times relative to Shared; 1,818 of 2,736 shaped model responses exceed the forty-frame minimum. The remote Projected/Full byte ratio is .932, with paired interval [.868,1.003]. These costs and the unequal recognition changes distinguish a measured shaping policy from a confidentiality guarantee.

Independent offline checks reconstruct all packet-derived aggregate and sequence features, fixture outcomes, LR selections and recognition arithmetic, event matching and cost accounting. Saved LR/forest inference is replayed with separate array arithmetic; CNN replay uses its exported tensors and the same numerical library. The package records the checking scope, implementation corrections and private raw-capture verification separately.

\subsection{Saved prefix and execution-path diagnostics}
Table~\ref{tab:prefix-diagnostics} exposes the prefix permutation results already computed in the original analysis. Qwen Full LR's 25/120 has unadjusted $p=.0595$; Llama's 19/120 has $p=.6092$. The within-session permutation preserves each session's label composition and fixes the predictor. These diagnostics neither demonstrate an information-free prefix nor change the primary post--prefix family.

Table~\ref{tab:paired-path-diagnostics} audits existing execution records, pairing held-out tasks by executor, session and unit. Each Full/Projected local arm has 518 executed tool actions, and all 120 pairs have equal call counts. Each arm also has 26 extra prerequisite-check calls. Exact action sequences agree in 85 Qwen and 94 Llama pairs; unordered action multisets agree in 118 and 119 pairs. Applied business actions agree in the same respective counts. Hence order changes and a few changed decisions remain even though call depth and repeated-check totals are matched. These summaries cannot identify the contribution of message contents, encoded size, action order or response timing to the recognition contrast.

\paragraph{What a separating design must preserve}
A future factorial contrast would need task-valid supplementary-content replacement and neutral padding matched at the observed encoding layer. Changing a record's category can change its semantics; matching JSON bytes need not match token counts, responses or timing. Record-count sweeps must account for their simultaneous size and workload changes. Public tool schemas could identify candidate fields, while prevalence of realized leakage would require sampled executions and an explicit denominator.

A new confirmation cohort would fix its estimand, validation-based attacker selection and multiplicity family before collection. Planning should vary plausible effects and session dependence; the pilot's fitted predictions alone do not establish the required sample size. Pooling models would define a weighted, model-stratified effect and retain each stratum's result. The existing LR/k-FP tests and exploratory RF/CNN results keep their original roles.

\begin{table*}[t]\centering\small
\caption{Saved prefix diagnostics. These are the original 2,000 within-session label permutations on fixed predictions, without multiplicity adjustment. They test association conditional on the observed sessions; a large value does not establish an information-free prefix.}\label{tab:prefix-diagnostics}
\setlength{\tabcolsep}{3pt}
\begin{tabularx}{\textwidth}{@{}lYrrrrr@{}}\toprule
Executor & Policy & $n$ & LR correct & LR $p$ & k-FP correct & k-FP $p$ \\\midrule
Qwen3:8B & Full & 120 & 25 & 0.0595 & 18 & 0.8071 \\
Qwen3:8B & Projected & 120 & 23 & 0.3998 & 23 & 0.4473 \\
Llama3.1:8B & Full & 120 & 19 & 0.6092 & 14 & 0.8646 \\
Llama3.1:8B & Projected & 120 & 16 & 0.5347 & 20 & 0.3988 \\
Hosted alias & Full & 48 & 5 & 0.5412 & 3 & 0.9960 \\
Hosted alias & Projected & 48 & 5 & 0.8141 & 6 & 0.9150 \\
\bottomrule\end{tabularx}
\end{table*}

\begin{table*}[t]\centering\small
\caption{Descriptive audit of the 120 held-out Full/Projected task pairs per local executor. Calls are executed tool actions. Extra checks count repeated account/policy/approval/status calls beyond the first. The last four columns count agreeing pairs; a multiset ignores call order. These summaries do not assign the recognition change to a mediator.}\label{tab:paired-path-diagnostics}
\setlength{\tabcolsep}{3pt}
\begin{tabularx}{\textwidth}{@{}lYrrrrr@{}}\toprule
Executor & Total calls F/P & Extra checks F/P & Same count & Same sequence & Same multiset & Same applied action \\\midrule
Qwen3:8B & 518/518 & 26/26 & 120 & 85 & 118 & 118 \\
Llama3.1:8B & 518/518 & 26/26 & 120 & 94 & 119 & 119 \\
\bottomrule\end{tabularx}
\end{table*}


\begin{table*}[t]\centering\small
\caption{Condition-matched recognition: correct predictions out of $n$ (six-class reference $n/6$). Prefix is the initial Full model call; Post is its first call after the ticket. Complete spans all requests and responses. RF and CNN are exploratory (E); neither replaces LR/k-FP in the primary family. $p_{proj}$ is the exact Holm-adjusted Full--Projected Post test. Complete-view comparisons are descriptive. Supplement S1 reports all uncertainty and frozen-transfer results.}\label{tab:phase-complete}
\setlength{\tabcolsep}{3pt}
\begin{tabularx}{\textwidth}{@{}lYrrrrrrrrr@{}}\toprule
Executor & Attack & $n$ & Prefix & \multicolumn{2}{c}{Post M} & \multicolumn{2}{c}{Complete M} & \multicolumn{2}{c}{Complete T+M} & $p_{proj}$ \\ & & & & Full & Proj. & Full & Proj. & Full & Proj. & \\\midrule
Qwen3:8B & LR & 120 & 25 & 43 & 23 & 46 & 44 & 54 & 39 & 0.0778 \\
Qwen3:8B & kFP & 120 & 18 & 32 & 20 & 43 & 39 & 39 & 30 & 0.1989 \\
Qwen3:8B & RF (E) & 120 & 25 & 50 & 30 & 61 & 31 & 55 & 30 & --- \\
Qwen3:8B & CNN (E) & 120 & 20 & 20 & 20 & 22 & 22 & 20 & 22 & --- \\
Llama3.1:8B & LR & 120 & 19 & 40 & 22 & 62 & 33 & 53 & 29 & 0.1172 \\
Llama3.1:8B & kFP & 120 & 14 & 31 & 20 & 47 & 34 & 45 & 31 & 0.1989 \\
Llama3.1:8B & RF (E) & 120 & 24 & 59 & 19 & 73 & 34 & 55 & 30 & --- \\
Llama3.1:8B & CNN (E) & 120 & 20 & 26 & 20 & 22 & 15 & 19 & 18 & --- \\
Hosted alias & LR & 48 & 5 & 15 & 10 & 9 & 9 & 12 & 6 & 0.9214 \\
Hosted alias & kFP & 48 & 3 & 6 & 7 & 9 & 5 & 13 & 11 & 1.0000 \\
Hosted alias & RF (E) & 48 & 8 & 15 & 9 & 7 & 7 & 13 & 8 & --- \\
Hosted alias & CNN (E) & 48 & 8 & 11 & 5 & 10 & 8 & 9 & 11 & --- \\
\bottomrule\end{tabularx}
\end{table*}

\begin{table*}[t]\centering\small
\caption{Paired feedback-path results. All scores are correct intended-action predictions in the model-link view. Models are fitted within each condition. Full minus Projected is a percentage-point difference; LR intervals resample training and test sessions with refitting, while k-FP intervals condition on fitted predictors. The adjusted $p$ columns enumerate the session-sign tests on saved paired predictions exactly; eight local tests and four remote tests form separate Holm families. The declared Monte Carlo approximation is retained beside the exact values in the supplement. Hosted results use a separate cohort and mutable alias.}\label{tab:randomized-primary}
\setlength{\tabcolsep}{4pt}
\begin{tabularx}{\textwidth}{@{}llrrrYrr@{}}\toprule
Executor & Attack & Prefix & \shortstack{Full\\post} & \shortstack{Projected\\post} & Full minus Projected [95\% CI] & $p_{post}$ & $p_{proj}$ \\\midrule
Qwen3:8B & LR & 25/120 & 43/120 & 23/120 & $+16.7\;[+4.0,+29.8]$ & 0.1357 & 0.0778 \\
Qwen3:8B & kFP & 18/120 & 32/120 & 20/120 & $+10.0\;[-1.6,+20.5]$ & 0.1172 & 0.1989 \\
Llama3.1:8B & LR & 19/120 & 40/120 & 22/120 & $+15.0\;[+3.1,+32.8]$ & 0.0508 & 0.1172 \\
Llama3.1:8B & kFP & 14/120 & 31/120 & 20/120 & $+9.2\;[+0.0,+17.4]$ & 0.1357 & 0.1989 \\
Hosted alias & LR & 5/48 & 15/48 & 10/48 & $+10.4\;[-3.2,+41.5]$ & 0.1562 & 0.9214 \\
Hosted alias & kFP & 3/48 & 6/48 & 7/48 & $-2.1\;[-17.4,+15.4]$ & 1.0000 & 1.0000 \\
\bottomrule\end{tabularx}
\end{table*}

\begin{table*}[t]\centering\small
\caption{Event overlap sensitivity on the same predictions. Each condition contains 120 tasks in 24 sessions, including four empty sessions. Learned detectors use T+M; all-active emits one session interval per class. IoU .5 is primary; .3/.7 are sensitivity checks. FP/session divides all false events by 24, retaining empty sessions. Supplement S1 retains hourly normalization, counts and intervals for these bounded recordings.}\label{tab:randomized-events}
\setlength{\tabcolsep}{3pt}
\begin{tabularx}{\textwidth}{@{}llYrrrrrr@{}}\toprule
Executor & Policy & Detector & F1 (.3) & F1 (.5) & F1 (.7) & P (.5) & R (.5) & FP/session (.5) \\\midrule
Qwen3:8B & Shared & LR & 0.216 & 0.119 & 0.048 & 0.072 & 0.350 & 22.54 \\
Qwen3:8B & Shared & RF & 0.106 & 0.077 & 0.035 & 0.043 & 0.400 & 44.79 \\
Qwen3:8B & Shared & all-active & 0.295 & 0.038 & 0.000 & 0.035 & 0.042 & 5.79 \\
Qwen3:8B & Joint & LR & 0.196 & 0.117 & 0.050 & 0.070 & 0.375 & 25.04 \\
Qwen3:8B & Joint & RF & 0.329 & 0.188 & 0.090 & 0.123 & 0.400 & 14.29 \\
Qwen3:8B & Joint & all-active & 0.265 & 0.038 & 0.000 & 0.035 & 0.042 & 5.79 \\
Llama3.1:8B & Shared & LR & 0.333 & 0.216 & 0.104 & 0.144 & 0.433 & 12.88 \\
Llama3.1:8B & Shared & RF & 0.464 & 0.249 & 0.105 & 0.186 & 0.375 & 8.21 \\
Llama3.1:8B & Shared & all-active & 0.508 & 0.098 & 0.000 & 0.090 & 0.108 & 5.46 \\
Llama3.1:8B & Joint & LR & 0.441 & 0.251 & 0.104 & 0.186 & 0.383 & 8.38 \\
Llama3.1:8B & Joint & RF & 0.429 & 0.274 & 0.129 & 0.199 & 0.442 & 8.92 \\
Llama3.1:8B & Joint & all-active & 0.530 & 0.121 & 0.000 & 0.111 & 0.133 & 5.33 \\
\bottomrule\end{tabularx}
\end{table*}

\begin{table*}[t]\centering\small
\caption{Numerical precision sensitivity for the complete declared primary families. MC retains the prespecified 10,000-draw session-sign approximation; exact enumerates the same sign-flip distribution using integer counts. This post-scoring arithmetic audit changes no predictions, clusters, hypotheses or multiplicity families. A PASS audit means arithmetic consistency, not hypothesis rejection.}\label{tab:randomized-sign-precision}
\setlength{\tabcolsep}{4pt}
\begin{tabularx}{\textwidth}{@{}llYrrrr@{}}\toprule
Executor & Attack & Contrast & MC $p$ & Exact $p$ & MC Holm & Exact Holm \\\midrule
Llama3.1:8B & LR & Post--prefix & 0.006199 & 0.006348 & 0.049595 & 0.050781 \\
Llama3.1:8B & kFP & Post--prefix & 0.031997 & 0.033936 & 0.127987 & 0.135742 \\
Llama3.1:8B & LR & Full--Projected & 0.021198 & 0.021362 & 0.109189 & 0.117188 \\
Llama3.1:8B & kFP & Full--Projected & 0.096890 & 0.099457 & 0.193781 & 0.198914 \\
Qwen3:8B & LR & Post--prefix & 0.036796 & 0.035339 & 0.127987 & 0.135742 \\
Qwen3:8B & kFP & Post--prefix & 0.018198 & 0.019531 & 0.109189 & 0.117188 \\
Qwen3:8B & LR & Full--Projected & 0.010499 & 0.011108 & 0.073493 & 0.077759 \\
Qwen3:8B & kFP & Full--Projected & 0.141986 & 0.144592 & 0.193781 & 0.198914 \\
Hosted alias & LR & Post--prefix & 0.039096 & 0.039062 & 0.156384 & 0.156250 \\
Hosted alias & kFP & Post--prefix & 0.558044 & 0.562500 & 1.000000 & 1.000000 \\
Hosted alias & LR & Full--Projected & 0.308369 & 0.307129 & 0.925107 & 0.921387 \\
Hosted alias & kFP & Full--Projected & 1.000000 & 1.000000 & 1.000000 & 1.000000 \\
\bottomrule\end{tabularx}
\end{table*}

\begin{table*}[t]\centering\small
\caption{Business completion and resource accounting. All-run and held-out success use distinct denominators. Remaining columns describe the held-out population. Median and P95 task times start at worker admission and exclude worker queueing. Session seconds/task amortize complete session time, including setup, background and idle margins; they are not per-task response latencies. Bytes count each packet once per session. Local and hosted absolute costs belong to different systems.}\label{tab:randomized-costs}
\setlength{\tabcolsep}{4pt}
\begin{tabularx}{\textwidth}{@{}lYrrrrrrr@{}}\toprule
Executor & Condition & \shortstack{All\\success} & \shortstack{Test\\success} & \shortstack{Median\\s} & \shortstack{P95\\s} & \shortstack{Session\\s/task} & \shortstack{T\\KiB/task} & \shortstack{M\\KiB/task} \\\midrule
Qwen3:8B & Full & 311/312 & 119/120 & 3.20 & 8.06 & 7.14 & 152.4 & 13.1 \\
Qwen3:8B & Projected & 311/312 & 119/120 & 3.04 & 7.89 & 6.91 & 151.4 & 10.2 \\
Qwen3:8B & Dedicated & 311/312 & 119/120 & 11.01 & 23.40 & 6.76 & 173.8 & 14.2 \\
Qwen3:8B & Shared & 311/312 & 119/120 & 9.61 & 31.78 & 6.91 & 144.6 & 12.7 \\
Qwen3:8B & Joint & 311/312 & 119/120 & 9.22 & 26.73 & 6.82 & 149.0 & 247.6 \\
Llama3.1:8B & Full & 309/312 & 117/120 & 6.84 & 9.62 & 9.82 & 165.4 & 13.1 \\
Llama3.1:8B & Projected & 306/312 & 118/120 & 6.82 & 9.74 & 9.94 & 167.4 & 10.3 \\
Llama3.1:8B & Dedicated & 309/312 & 117/120 & 21.38 & 29.28 & 9.22 & 179.7 & 14.2 \\
Llama3.1:8B & Shared & 309/312 & 117/120 & 20.94 & 28.99 & 9.04 & 151.3 & 12.7 \\
Llama3.1:8B & Joint & 309/312 & 117/120 & 23.07 & 32.50 & 9.42 & 161.0 & 286.2 \\
Hosted alias & Full & 92/96 & 46/48 & 14.49 & 33.24 & 20.76 & 256.1 & 108.5 \\
Hosted alias & Projected & 95/96 & 47/48 & 13.83 & 28.93 & 20.42 & 253.6 & 101.1 \\
\bottomrule\end{tabularx}
\end{table*}

\FloatBarrier
\clearpage\onecolumn
\small\setlength{\tabcolsep}{3pt}
\begin{longtable}{@{}P{.14\textwidth}P{.14\textwidth}P{.08\textwidth}P{.04\textwidth}P{.125\textwidth}P{.125\textwidth}P{.125\textwidth}P{.125\textwidth}@{}}
\caption{Complete recognition outcomes. Parentheses give the test denominator. Each attack cell gives correct predictions as matched fit / frozen Full fit; a dash means the same Full model is already reported. Prefix and post-ticket use M; complete-execution views are separately fitted. Every planned terminal outcome remains in the intended-label denominator.}\label{tab:randomized-all-recognition}\\
\toprule
Executor ($n$) & Condition & Phase & View & LR & k-FP & RF & CNN \\\midrule\endfirsthead
\multicolumn{8}{l}{\textit{Table \thetable{} continued}}\\
\toprule
Executor ($n$) & Condition & Phase & View & LR & k-FP & RF & CNN \\\midrule\endhead
\bottomrule\endfoot
Qwen3:8B (120) & Full & Prefix & M & 25 / -- & 18 / -- & 25 / -- & 20 / -- \\
Qwen3:8B (120) & Full & Post-ticket & M & 43 / -- & 32 / -- & 50 / -- & 20 / -- \\
Qwen3:8B (120) & Full & Complete & T & 19 / -- & 21 / -- & 23 / -- & 20 / -- \\
Qwen3:8B (120) & Full & Complete & M & 46 / -- & 43 / -- & 61 / -- & 22 / -- \\
Qwen3:8B (120) & Full & Complete & T+M & 54 / -- & 39 / -- & 55 / -- & 20 / -- \\
Qwen3:8B (120) & Projected & Prefix & M & 23 / 24 & 23 / 19 & 21 / 16 & 20 / 20 \\
Qwen3:8B (120) & Projected & Post-ticket & M & 23 / 20 & 20 / 20 & 30 / 20 & 20 / 20 \\
Qwen3:8B (120) & Projected & Complete & T & 18 / 20 & 20 / 24 & 20 / 24 & 20 / 22 \\
Qwen3:8B (120) & Projected & Complete & M & 44 / 20 & 39 / 30 & 31 / 20 & 22 / 16 \\
Qwen3:8B (120) & Projected & Complete & T+M & 39 / 19 & 30 / 28 & 30 / 16 & 22 / 22 \\
Qwen3:8B (120) & Dedicated & Complete & T & 23 / 21 & 25 / 20 & 24 / 20 & 20 / 20 \\
Qwen3:8B (120) & Dedicated & Complete & M & 25 / 24 & 22 / 19 & 24 / 23 & 25 / 19 \\
Qwen3:8B (120) & Dedicated & Complete & T+M & 30 / 28 & 23 / 21 & 24 / 20 & 16 / 21 \\
Qwen3:8B (120) & Shared & Complete & T & 20 / 18 & 19 / 21 & 17 / 19 & 20 / 20 \\
Qwen3:8B (120) & Shared & Complete & M & 25 / 26 & 25 / 19 & 30 / 24 & 17 / 20 \\
Qwen3:8B (120) & Shared & Complete & T+M & 23 / 21 & 22 / 20 & 22 / 24 & 19 / 19 \\
Qwen3:8B (120) & Joint & Complete & T & 19 / 24 & 17 / 16 & 19 / 19 & 17 / 19 \\
Qwen3:8B (120) & Joint & Complete & M & 18 / 20 & 24 / 20 & 22 / 20 & 25 / 20 \\
Qwen3:8B (120) & Joint & Complete & T+M & 20 / 20 & 22 / 21 & 19 / 19 & 25 / 21 \\
Llama3.1:8B (120) & Full & Prefix & M & 19 / -- & 14 / -- & 24 / -- & 20 / -- \\
Llama3.1:8B (120) & Full & Post-ticket & M & 40 / -- & 31 / -- & 59 / -- & 26 / -- \\
Llama3.1:8B (120) & Full & Complete & T & 17 / -- & 18 / -- & 19 / -- & 17 / -- \\
Llama3.1:8B (120) & Full & Complete & M & 62 / -- & 47 / -- & 73 / -- & 22 / -- \\
Llama3.1:8B (120) & Full & Complete & T+M & 53 / -- & 45 / -- & 55 / -- & 19 / -- \\
Llama3.1:8B (120) & Projected & Prefix & M & 16 / 20 & 20 / 12 & 14 / 21 & 20 / 20 \\
Llama3.1:8B (120) & Projected & Post-ticket & M & 22 / 20 & 20 / 27 & 19 / 20 & 20 / 21 \\
Llama3.1:8B (120) & Projected & Complete & T & 18 / 22 & 19 / 20 & 16 / 23 & 19 / 17 \\
Llama3.1:8B (120) & Projected & Complete & M & 33 / 20 & 34 / 23 & 34 / 16 & 15 / 17 \\
Llama3.1:8B (120) & Projected & Complete & T+M & 29 / 20 & 31 / 27 & 30 / 19 & 18 / 20 \\
Llama3.1:8B (120) & Dedicated & Complete & T & 16 / 21 & 17 / 19 & 15 / 20 & 21 / 21 \\
Llama3.1:8B (120) & Dedicated & Complete & M & 17 / 24 & 13 / 16 & 13 / 20 & 19 / 18 \\
Llama3.1:8B (120) & Dedicated & Complete & T+M & 20 / 23 & 16 / 21 & 18 / 20 & 19 / 18 \\
Llama3.1:8B (120) & Shared & Complete & T & 26 / 19 & 13 / 15 & 22 / 19 & 16 / 20 \\
Llama3.1:8B (120) & Shared & Complete & M & 19 / 21 & 19 / 17 & 21 / 21 & 19 / 18 \\
Llama3.1:8B (120) & Shared & Complete & T+M & 19 / 22 & 17 / 17 & 17 / 21 & 20 / 20 \\
Llama3.1:8B (120) & Joint & Complete & T & 17 / 18 & 23 / 17 & 15 / 22 & 24 / 19 \\
Llama3.1:8B (120) & Joint & Complete & M & 25 / 20 & 22 / 20 & 16 / 20 & 21 / 20 \\
Llama3.1:8B (120) & Joint & Complete & T+M & 22 / 20 & 17 / 20 & 17 / 20 & 18 / 18 \\
Hosted alias (48) & Full & Prefix & M & 5 / -- & 3 / -- & 8 / -- & 8 / -- \\
Hosted alias (48) & Full & Post-ticket & M & 15 / -- & 6 / -- & 15 / -- & 11 / -- \\
Hosted alias (48) & Full & Complete & T & 8 / -- & 8 / -- & 10 / -- & 11 / -- \\
Hosted alias (48) & Full & Complete & M & 9 / -- & 9 / -- & 7 / -- & 10 / -- \\
Hosted alias (48) & Full & Complete & T+M & 12 / -- & 13 / -- & 13 / -- & 9 / -- \\
Hosted alias (48) & Projected & Prefix & M & 5 / 5 & 6 / 5 & 6 / 8 & 8 / 8 \\
Hosted alias (48) & Projected & Post-ticket & M & 10 / 7 & 7 / 6 & 9 / 7 & 5 / 11 \\
Hosted alias (48) & Projected & Complete & T & 5 / 8 & 10 / 7 & 8 / 8 & 7 / 13 \\
Hosted alias (48) & Projected & Complete & M & 9 / 7 & 5 / 10 & 7 / 6 & 8 / 11 \\
Hosted alias (48) & Projected & Complete & T+M & 6 / 6 & 11 / 8 & 8 / 8 & 11 / 13 \\
\end{longtable}
\normalsize\clearpage\twocolumn

\clearpage\onecolumn
\small\setlength{\tabcolsep}{3pt}
\begin{longtable}{@{}P{.14\textwidth}P{.13\textwidth}P{.05\textwidth}P{.21\textwidth}rrrrr@{}}
\caption{Complete event results at the primary IoU .5. P, R and F1 use one-to-one task-instance matching. Empty FP counts false events in empty test sessions. The artifact additionally provides .3/.7 overlap results, validation choices, per-session counts and uncertainty.}\label{tab:randomized-all-events}\\
\toprule
Executor & Condition & View & Detector & P & R & F1 & FP/hour & Empty FP \\\midrule\endfirsthead
\multicolumn{9}{l}{\textit{Table \thetable{} continued}}\\
\toprule
Executor & Condition & View & Detector & P & R & F1 & FP/hour & Empty FP \\\midrule\endhead
\bottomrule\endfoot
Qwen3:8B & Dedicated & T & LR & 0.080 & 0.292 & 0.125 & 1788.5 & 15 \\
Qwen3:8B & Dedicated & T & RF & 0.051 & 0.317 & 0.088 & 3119.8 & 70 \\
Qwen3:8B & Dedicated & T & gap & 0.060 & 0.058 & 0.059 & 483.7 & 16 \\
Qwen3:8B & Dedicated & M & LR & 0.028 & 0.250 & 0.051 & 4606.5 & 3 \\
Qwen3:8B & Dedicated & M & RF & 0.079 & 0.317 & 0.127 & 1961.5 & 6 \\
Qwen3:8B & Dedicated & M & gap & 0.041 & 0.075 & 0.053 & 923.1 & 4 \\
Qwen3:8B & Dedicated & T+M & LR & 0.033 & 0.300 & 0.059 & 4726.3 & 3 \\
Qwen3:8B & Dedicated & T+M & RF & 0.139 & 0.392 & 0.205 & 1291.4 & 0 \\
Qwen3:8B & Dedicated & T+M & gap & 0.056 & 0.042 & 0.048 & 372.8 & 16 \\
Qwen3:8B & Full & T & LR & 0.011 & 0.033 & 0.016 & 1563.4 & 2 \\
Qwen3:8B & Full & T & RF & 0.009 & 0.142 & 0.017 & 7673.9 & 60 \\
Qwen3:8B & Full & T & gap & 0.006 & 0.025 & 0.010 & 2118.1 & 30 \\
Qwen3:8B & Full & M & LR & 0.054 & 0.358 & 0.094 & 3168.7 & 0 \\
Qwen3:8B & Full & M & RF & 0.022 & 0.317 & 0.041 & 7136.0 & 0 \\
Qwen3:8B & Full & M & gap & 0.005 & 0.008 & 0.006 & 882.5 & 4 \\
Qwen3:8B & Full & T+M & LR & 0.049 & 0.383 & 0.087 & 3748.7 & 0 \\
Qwen3:8B & Full & T+M & RF & 0.036 & 0.408 & 0.067 & 5471.7 & 0 \\
Qwen3:8B & Full & T+M & gap & 0.011 & 0.025 & 0.015 & 1159.9 & 23 \\
Qwen3:8B & Projected & T & LR & 0.007 & 0.067 & 0.013 & 4986.1 & 56 \\
Qwen3:8B & Projected & T & RF & 0.010 & 0.117 & 0.018 & 6019.9 & 40 \\
Qwen3:8B & Projected & T & gap & 0.000 & 0.000 & 0.000 & 1120.6 & 23 \\
Qwen3:8B & Projected & M & LR & 0.047 & 0.208 & 0.077 & 2189.0 & 0 \\
Qwen3:8B & Projected & M & RF & 0.003 & 0.033 & 0.006 & 5355.3 & 0 \\
Qwen3:8B & Projected & M & gap & 0.015 & 0.025 & 0.019 & 833.9 & 4 \\
Qwen3:8B & Projected & T+M & LR & 0.045 & 0.300 & 0.079 & 3283.6 & 0 \\
Qwen3:8B & Projected & T+M & RF & 0.020 & 0.192 & 0.036 & 4881.9 & 2 \\
Qwen3:8B & Projected & T+M & gap & 0.008 & 0.017 & 0.010 & 1133.6 & 23 \\
Qwen3:8B & Shared & T & LR & 0.037 & 0.217 & 0.064 & 2901.0 & 86 \\
Qwen3:8B & Shared & T & RF & 0.049 & 0.367 & 0.087 & 3704.4 & 77 \\
Qwen3:8B & Shared & T & gap & 0.073 & 0.025 & 0.037 & 165.0 & 9 \\
Qwen3:8B & Shared & M & LR & 0.043 & 0.367 & 0.077 & 4234.2 & 16 \\
Qwen3:8B & Shared & M & RF & 0.080 & 0.375 & 0.132 & 2245.2 & 0 \\
Qwen3:8B & Shared & M & gap & 0.056 & 0.092 & 0.069 & 812.1 & 4 \\
Qwen3:8B & Shared & T+M & LR & 0.072 & 0.350 & 0.119 & 2349.4 & 21 \\
Qwen3:8B & Shared & T+M & RF & 0.043 & 0.400 & 0.077 & 4668.5 & 2 \\
Qwen3:8B & Shared & T+M & gap & 0.023 & 0.050 & 0.031 & 1129.1 & 23 \\
Qwen3:8B & Joint & T & LR & 0.022 & 0.300 & 0.041 & 6978.7 & 138 \\
Qwen3:8B & Joint & T & RF & 0.037 & 0.333 & 0.067 & 4555.7 & 82 \\
Qwen3:8B & Joint & T & gap & 0.034 & 0.025 & 0.029 & 373.8 & 16 \\
Qwen3:8B & Joint & M & LR & 0.121 & 0.375 & 0.183 & 1433.6 & 7 \\
Qwen3:8B & Joint & M & RF & 0.092 & 0.408 & 0.151 & 2115.2 & 0 \\
Qwen3:8B & Joint & M & gap & 0.049 & 0.017 & 0.025 & 171.5 & 4 \\
Qwen3:8B & Joint & T+M & LR & 0.070 & 0.375 & 0.117 & 2642.8 & 0 \\
Qwen3:8B & Joint & T+M & RF & 0.123 & 0.400 & 0.188 & 1508.3 & 3 \\
Qwen3:8B & Joint & T+M & gap & 0.015 & 0.033 & 0.021 & 1156.5 & 23 \\
Qwen3:8B & Dedicated & none & all-active & 0.028 & 0.033 & 0.030 & 621.3 & 24 \\
Qwen3:8B & Dedicated & none & training-presence-prior & 0.028 & 0.017 & 0.021 & 310.6 & 12 \\
Qwen3:8B & Full & none & all-active & 0.000 & 0.000 & 0.000 & 605.2 & 24 \\
Qwen3:8B & Full & none & training-presence-prior & 0.000 & 0.000 & 0.000 & 302.6 & 12 \\
Qwen3:8B & Projected & none & all-active & 0.000 & 0.000 & 0.000 & 625.4 & 24 \\
Qwen3:8B & Projected & none & training-presence-prior & 0.000 & 0.000 & 0.000 & 312.7 & 12 \\
Qwen3:8B & Shared & none & all-active & 0.035 & 0.042 & 0.038 & 603.6 & 24 \\
Qwen3:8B & Shared & none & training-presence-prior & 0.056 & 0.033 & 0.042 & 295.3 & 12 \\
Qwen3:8B & Joint & none & all-active & 0.035 & 0.042 & 0.038 & 611.2 & 24 \\
Qwen3:8B & Joint & none & training-presence-prior & 0.056 & 0.033 & 0.042 & 299.0 & 12 \\
\pagebreak
Llama3.1:8B & Dedicated & T & LR & 0.086 & 0.267 & 0.130 & 1105.9 & 61 \\
Llama3.1:8B & Dedicated & T & RF & 0.061 & 0.425 & 0.106 & 2569.6 & 83 \\
Llama3.1:8B & Dedicated & T & gap & 0.070 & 0.025 & 0.037 & 130.1 & 9 \\
Llama3.1:8B & Dedicated & M & LR & 0.194 & 0.358 & 0.251 & 582.2 & 2 \\
Llama3.1:8B & Dedicated & M & RF & 0.140 & 0.450 & 0.213 & 1079.9 & 0 \\
Llama3.1:8B & Dedicated & M & gap & 0.143 & 0.050 & 0.074 & 117.1 & 4 \\
Llama3.1:8B & Dedicated & T+M & LR & 0.169 & 0.425 & 0.242 & 813.2 & 1 \\
Llama3.1:8B & Dedicated & T+M & RF & 0.144 & 0.400 & 0.212 & 927.0 & 0 \\
Llama3.1:8B & Dedicated & T+M & gap & 0.055 & 0.050 & 0.052 & 338.3 & 16 \\
Llama3.1:8B & Full & T & LR & 0.027 & 0.158 & 0.046 & 2080.6 & 60 \\
Llama3.1:8B & Full & T & RF & 0.008 & 0.142 & 0.015 & 6593.2 & 56 \\
Llama3.1:8B & Full & T & gap & 0.000 & 0.000 & 0.000 & 2083.7 & 33 \\
Llama3.1:8B & Full & M & LR & 0.082 & 0.650 & 0.145 & 2676.4 & 0 \\
Llama3.1:8B & Full & M & RF & 0.035 & 0.492 & 0.066 & 4906.7 & 0 \\
Llama3.1:8B & Full & M & gap & 0.000 & 0.000 & 0.000 & 128.3 & 4 \\
Llama3.1:8B & Full & T+M & LR & 0.073 & 0.633 & 0.130 & 2969.7 & 0 \\
Llama3.1:8B & Full & T+M & RF & 0.037 & 0.492 & 0.069 & 4680.6 & 1 \\
Llama3.1:8B & Full & T+M & gap & 0.013 & 0.008 & 0.010 & 226.1 & 16 \\
Llama3.1:8B & Projected & T & LR & 0.022 & 0.142 & 0.038 & 2316.9 & 67 \\
Llama3.1:8B & Projected & T & RF & 0.007 & 0.125 & 0.013 & 6718.4 & 51 \\
Llama3.1:8B & Projected & T & gap & 0.000 & 0.000 & 0.000 & 2084.6 & 33 \\
Llama3.1:8B & Projected & M & LR & 0.040 & 0.142 & 0.062 & 1245.9 & 0 \\
Llama3.1:8B & Projected & M & RF & 0.009 & 0.167 & 0.017 & 6600.8 & 0 \\
Llama3.1:8B & Projected & M & gap & 0.000 & 0.000 & 0.000 & 1695.4 & 4 \\
Llama3.1:8B & Projected & T+M & LR & 0.037 & 0.242 & 0.065 & 2250.5 & 0 \\
Llama3.1:8B & Projected & T+M & RF & 0.023 & 0.350 & 0.044 & 5285.4 & 0 \\
Llama3.1:8B & Projected & T+M & gap & 0.000 & 0.000 & 0.000 & 2081.6 & 33 \\
Llama3.1:8B & Shared & T & LR & 0.114 & 0.350 & 0.171 & 1088.6 & 63 \\
Llama3.1:8B & Shared & T & RF & 0.059 & 0.433 & 0.103 & 2774.5 & 91 \\
Llama3.1:8B & Shared & T & gap & 0.071 & 0.083 & 0.077 & 431.4 & 16 \\
Llama3.1:8B & Shared & M & LR & 0.145 & 0.392 & 0.212 & 919.3 & 20 \\
Llama3.1:8B & Shared & M & RF & 0.174 & 0.508 & 0.259 & 962.5 & 0 \\
Llama3.1:8B & Shared & M & gap & 0.102 & 0.042 & 0.059 & 146.0 & 4 \\
Llama3.1:8B & Shared & T+M & LR & 0.144 & 0.433 & 0.216 & 1025.5 & 17 \\
Llama3.1:8B & Shared & T+M & RF & 0.186 & 0.375 & 0.249 & 653.8 & 0 \\
Llama3.1:8B & Shared & T+M & gap & 0.059 & 0.017 & 0.026 & 106.2 & 9 \\
Llama3.1:8B & Joint & T & LR & 0.107 & 0.250 & 0.150 & 795.8 & 47 \\
Llama3.1:8B & Joint & T & RF & 0.066 & 0.483 & 0.117 & 2594.2 & 143 \\
Llama3.1:8B & Joint & T & gap & 0.079 & 0.058 & 0.067 & 261.0 & 16 \\
Llama3.1:8B & Joint & M & LR & 0.182 & 0.358 & 0.242 & 614.3 & 4 \\
Llama3.1:8B & Joint & M & RF & 0.175 & 0.467 & 0.255 & 840.3 & 0 \\
Llama3.1:8B & Joint & M & gap & 0.205 & 0.075 & 0.110 & 111.4 & 4 \\
Llama3.1:8B & Joint & T+M & LR & 0.186 & 0.383 & 0.251 & 639.8 & 2 \\
Llama3.1:8B & Joint & T+M & RF & 0.199 & 0.442 & 0.274 & 681.2 & 2 \\
Llama3.1:8B & Joint & T+M & gap & 0.129 & 0.075 & 0.095 & 194.2 & 16 \\
Llama3.1:8B & Dedicated & none & all-active & 0.083 & 0.100 & 0.091 & 429.3 & 24 \\
Llama3.1:8B & Dedicated & none & training-presence-prior & 0.069 & 0.042 & 0.052 & 217.9 & 12 \\
Llama3.1:8B & Full & none & all-active & 0.000 & 0.000 & 0.000 & 440.0 & 24 \\
Llama3.1:8B & Full & none & training-presence-prior & 0.000 & 0.000 & 0.000 & 220.0 & 12 \\
Llama3.1:8B & Projected & none & all-active & 0.000 & 0.000 & 0.000 & 434.4 & 24 \\
Llama3.1:8B & Projected & none & training-presence-prior & 0.000 & 0.000 & 0.000 & 217.2 & 12 \\
Llama3.1:8B & Shared & none & all-active & 0.090 & 0.108 & 0.098 & 434.8 & 24 \\
Llama3.1:8B & Shared & none & training-presence-prior & 0.069 & 0.042 & 0.052 & 222.4 & 12 \\
Llama3.1:8B & Joint & none & all-active & 0.111 & 0.133 & 0.121 & 407.4 & 24 \\
Llama3.1:8B & Joint & none & training-presence-prior & 0.111 & 0.067 & 0.083 & 203.7 & 12 \\
\pagebreak
Hosted alias & Full & T & LR & 0.023 & 0.292 & 0.042 & 2193.0 & 21 \\
Hosted alias & Full & T & RF & 0.001 & 0.021 & 0.001 & 4895.3 & 28 \\
Hosted alias & Full & T & gap & 0.000 & 0.000 & 0.000 & 751.5 & 12 \\
Hosted alias & Full & M & LR & 0.024 & 0.458 & 0.047 & 3164.8 & 1 \\
Hosted alias & Full & M & RF & 0.008 & 0.208 & 0.015 & 4497.9 & 0 \\
Hosted alias & Full & M & gap & 0.008 & 0.042 & 0.013 & 928.5 & 1 \\
Hosted alias & Full & T+M & LR & 0.016 & 0.375 & 0.031 & 3974.1 & 1 \\
Hosted alias & Full & T+M & RF & 0.011 & 0.292 & 0.021 & 4519.6 & 0 \\
Hosted alias & Full & T+M & gap & 0.009 & 0.062 & 0.016 & 1156.1 & 13 \\
Hosted alias & Projected & T & LR & 0.057 & 0.479 & 0.101 & 1406.8 & 12 \\
Hosted alias & Projected & T & RF & 0.000 & 0.000 & 0.000 & 0.0 & 0 \\
Hosted alias & Projected & T & gap & 0.000 & 0.000 & 0.000 & 1792.5 & 13 \\
Hosted alias & Projected & M & LR & 0.027 & 0.354 & 0.050 & 2244.3 & 0 \\
Hosted alias & Projected & M & RF & 0.002 & 0.062 & 0.004 & 4969.9 & 0 \\
Hosted alias & Projected & M & gap & 0.062 & 0.042 & 0.050 & 110.2 & 1 \\
Hosted alias & Projected & T+M & LR & 0.038 & 0.521 & 0.071 & 2321.5 & 0 \\
Hosted alias & Projected & T+M & RF & 0.004 & 0.104 & 0.007 & 4712.7 & 1 \\
Hosted alias & Projected & T+M & gap & 0.019 & 0.042 & 0.026 & 374.7 & 13 \\
Hosted alias & Full & none & all-active & 0.009 & 0.021 & 0.013 & 386.6 & 12 \\
Hosted alias & Full & none & training-presence-prior & 0.014 & 0.021 & 0.017 & 256.5 & 8 \\
Hosted alias & Projected & none & all-active & 0.000 & 0.000 & 0.000 & 396.7 & 12 \\
Hosted alias & Projected & none & training-presence-prior & 0.000 & 0.000 & 0.000 & 264.5 & 8 \\
\end{longtable}
\normalsize\clearpage\twocolumn

